\Page{11}{\textenglish{03 / CAMPAIGNS}}{הדוגמנית שאולי לא קיימת}{43}{49}

{\fontsize{12.6}{18}\selectfont

\StoryColumns{בקמפיין קיץ, האור הוא חלק מהליהוק. הוא קובע אם הבגד ייראה קליל או כבד, ואם צבע שבלט בסטודיו עדיין יחזיק את העין בחוץ. בהפקה מסורתית, ההחלטות האלה נבחנות על הסט, כשצריך להחזיק יחד דוגמנית, סטיילינג, לוקיישן ומצלמה. את מה שמשתבש מנסים לתקן בעוד טייק, בהחלפת פריט או בעריכה. הרבה ממה שייראה טבעי בתמונה תלוי בעבודה שתוכננה מראש.}{המפיקה מתחילה מתפוקה. כמה תמונות צריך, האם מצלמים גם סרטון, והיכן כל גרסה תופיע? מנהלת האמנות והצלמת מחליטות איך הלוק יפגוש את המקום. בליהוק מחפשים תנועה, יציבה ונוכחות שמתאימות לבריף, ומסכמים מראש תנאי שימוש בתמונות. הסטייליסטית מזמינה דוגמאות ומתכננת גיבוי אם פריט לא יגיע במידה הרצויה. אנשי שיער, איפור ותאורה צריכים לעבוד באותו לוח זמנים. אם הלוקיישן בחוץ, שינוע הצוות והבגדים תלוי גם בגישה, במזג האוויר ובשעה שבה האור יגיע לכיוון הנכון.}

\par}

\Photo{fashion-show-backstage.jpg}{79mm}{מאחורי הקלעים של תצוגת אופנה, פברואר 2008. \BackstageCredit}

{\fontsize{12.2}{17.2}\selectfont

\StoryColumns{לקראת קמפיין של \textenglish{Puma} עם נועה קירל, מנהל השיווק פנה לדניאלה קרול בהתראה קצרה. לדבריה בראיון ל״פורטפוליו״, היא הכינה את הקריאייטיב בתוך שבוע והפיקה את יום הצילום מאמסטרדם, בתיאום עם הצוות בישראל. זה מקרה שבו ההפקה נוהלה מרחוק והצילום עצמו עדיין נשען על צוות בשטח. זמן ההכנה שתיארה קרול נוגע לקמפיין הזה בלבד. \SourceRef{7}}{הבחירה נמשכת גם אחרי שהסט מתקפל. עורכות וצוות המותג עוברות על החומר, בוחרות פריימים, מתקנות צבע ומתאימות את הקבצים לערוצים. לצד שכר הצוות, בתקציב עשויים להופיע מקום, ציוד, נסיעות, זכויות שימוש וזמן פוסט. כל אלה תלויים בבריף ובהיקף. יום צילום אחד יכול להניב תמונות לקטלוג ולרשתות, אך הוא עלול להצריך עבודת עריכה שונה מאוד לכל אחד מהשימושים האלה.}

\par}
